% Generated draft: requires mathematical and source review.
\documentclass[11pt]{article}
\usepackage[margin=25mm]{geometry}
\usepackage{fontspec}
\setmainfont{Latin Modern Roman}
\usepackage{amsmath,amssymb,mathtools,unicode-math}
\setmathfont{Latin Modern Math}
\usepackage{enumitem,needspace,xurl,hyperref}
\setlength{\parindent}{0pt}
\setlength{\parskip}{5pt}
\setlength{\emergencystretch}{3em}
\newcommand{\N}{\mathbb{N}}
\newcommand{\Z}{\mathbb{Z}}
\newcommand{\Q}{\mathbb{Q}}
\newcommand{\R}{\mathbb{R}}
\newcommand{\C}{\mathbb{C}}
\newcommand{\F}{\mathbb{F}}
\newcommand{\A}{\mathbb{A}}
\newcommand{\PP}{\mathbb{P}}
\newcommand{\E}{\mathbb{E}}
\newcommand{\eps}{\varepsilon}
\newcommand{\dd}{\,\mathrm{d}}
\begin{document}
\section*{UnsolvedMath ID: 30002044}
\textbf{Source problem code:} catalogue entry
\subsection{Definitions and mathematical statement}
The source asks for a useful geometric description of the generally nonconvex cone of \(q\)-ample line bundles on a variety. It does not specify the ground field, projectivity hypotheses, range of \(q\), ambient space of the cone, or the meaning of a useful description.

To explain the standard cohomological terminology in a conditional projective setting, let \(X\) be a projective variety over a field \(k\), and let \(q\geq 0\) be an integer. A line bundle \(L\) on \(X\) is \emph{\(q\)-ample} if, for every coherent sheaf \(\mathcal F\) on \(X\), there exists an integer \(m_0=m_0(\mathcal F,L)\geq 0\) such that
\[
 H^i\!\left(X,\mathcal F\otimes L^{\otimes m}\right)=0
 \qquad\text{for every integer }i>q\text{ and every integer }m\geq m_0.
\]
If \(q\geq\dim X\), this condition is automatic in this projective setting; the potentially nontrivial range is \(0\leq q<\dim X\). These projective hypotheses supply a conventional explanatory framework and are not stated hypotheses of the supplied question.

In this setting, numerical equivalence of line bundles means equality of their degrees on every integral curve in \(X\). Write
\[
 N^1(X)_{\R}
 =\bigl(\operatorname{Pic}(X)/{\equiv_{\mathrm{num}}}\bigr)
   \otimes_{\Z}\R,
 \qquad
 \mathcal S_q(X)=\{[L]_{\mathrm{num}}:L\text{ is a }q\text{-ample line bundle on }X\}.
\]
The numerical classes in \(\mathcal S_q(X)\) explain the line-bundle data underlying the customary numerical-cone interpretation. The supplied record does not specify how its cone is defined for real classes, so no particular extension from line bundles to real classes is imposed here. A generally nonconvex cone must not be replaced by the convex cone generated by its elements.

\textbf{Problem, retaining the source's unresolved scope.} Give a useful geometric description of the generally nonconvex cone of \(q\)-ample line bundles on a variety. Under a projective numerical interpretation, this asks for a geometric characterization of the \(q\)-ample locus in \(N^1(X)_{\R}\), with the real-class convention and hypotheses still to be specified. The source does not prescribe a concrete criterion or a precise standard for what counts as a useful description. Its toric background does not explicitly restrict the question to toric varieties.
\subsection{Short English statement}
A line bundle is \(q\)-ample when sufficiently high positive powers kill coherent-sheaf cohomology in every degree greater than \(q\). The question asks for a geometric description of the resulting generally nonconvex cone. The source leaves the class of varieties, the real-class convention, and the requested form of description unspecified.
\subsection{Sources}
\begin{itemize}
\item Local UnsolvedMath registry, permanent ID \(30002044\), title ``Geometry of Cones of \(q\)-Ample Line Bundles,'' source problem code \texttt{catalogue entry}, Oberwolfach Reports Open Problems collection. No individual proposer is supplied; the proposed year is recorded as 2012.
\item Supplied report citation: \emph{Toric Geometry (2013)}, Oberwolfach Reports, DOI \texttt{10.4171/owr/2012/21}, pp.\ 1231--1306. \url{https://doi.org/10.4171/owr/2012/21}. The citation's 2013 title date and the record's 2012 proposed year and DOI component are preserved without reconciliation.
\item Totaro, \emph{Line bundles with partially vanishing cohomology}. Supplied literature link: \url{https://ems.press/content/serial-article-files/31884?nt=1}.
\item Broomhead--Ottem--Prendergast-Smith, \emph{Partially ample line bundles on toric varieties}. Supplied literature link: \url{https://arxiv.org/abs/1202.3065}.
\item Literature assessment dated exactly \texttt{2026-08-21}: the supplied curation labels the problem open, reports substantial polyhedral descriptions for toric varieties, and attributes to Totaro the identification of a general numerical characterization of \(q\)-ampleness as open. It consequently regards the broad cone-description problem as open. This is supplied dated curation, not an independently verified literature assessment.
\item The registry labels the statement \texttt{corrected\_verified} and reports that the official report was checked, the complete question preserved, and its wording tightened without changing scope. These are supplied verification claims; no source consultation or mathematical verification is asserted by this draft.
\end{itemize}
\end{document}
